\documentclass[12pt, letterpaper]{article}
\usepackage[margin=1in]{geometry}
\usepackage{amsmath}
\usepackage{amssymb}
\usepackage{mathtools}
\usepackage{comment}
\usepackage{enumerate}
\usepackage{changepage}
\usepackage{amsthm}
\usepackage{bbm}
\usepackage{tikz-cd}
\usepackage{xcolor}
\usepackage{hyperref}

\setlength{\parindent}{0pt}

%% code from mathabx.sty and mathabx.dcl
\DeclareFontFamily{U}{mathx}{\hyphenchar\font45}
\DeclareFontShape{U}{mathx}{m}{n}{
      <5> <6> <7> <8> <9> <10>
      <10.95> <12> <14.4> <17.28> <20.74> <24.88>
      mathx10
      }{}
\DeclareSymbolFont{mathx}{U}{mathx}{m}{n}
\DeclareFontSubstitution{U}{mathx}{m}{n}
\DeclareMathAccent{\widecheck}{0}{mathx}{"71}
\DeclareMathAccent{\wideparen}{0}{mathx}{"75}

\def\cs#1{\texttt{\char`\\#1}}

\allowdisplaybreaks

\DeclareMathOperator{\Id}{Id}
\DeclareMathOperator{\im}{im}
\renewcommand{\Re}{\text{Re}}
\renewcommand{\Im}{\text{Im}}
\newcommand{\D}{\mathbf{D}}
\newcommand{\0}{\mathbf{0}}
\newcommand{\1}{\mathbf{1}}
\newcommand{\loc}{\text{loc}}
\DeclareMathOperator{\dist}{dist}
\DeclareMathOperator{\Span}{Span}
\DeclareMathOperator{\sign}{sign}
\DeclareMathOperator{\supp}{supp}
\DeclareMathOperator{\op}{op}
\DeclareMathOperator{\tr}{tr}
\DeclareMathOperator{\x}{\mathbf{x}}
\DeclareMathOperator{\y}{\mathbf{y}}
\DeclareMathOperator{\z}{\mathbf{z}}

\title{Math 104 Homework 6}
\author{Due Date: Monday, August 10 at 11:59pm}
\date{}

\begin{document}

\maketitle

\textbf{Instructions:} Submit pdf solutions in LaTeX (preferred), or \textit{neatly} handwritten. Unless otherwise stated, you may cite any result from class or from Rudin's book, but all other assertions you make must be proved rigorously. At the top of your submission, please write down the names of the other students you collaborated with and any outside sources you consulted.

\vspace{0.5cm}

\begin{enumerate}
    \item 
    Let $a \in \mathbb{R}$, and suppose $f : \mathbb{R} \to \mathbb{R}$ is Riemann integrable on $[a, b]$ for all $b > a$. We define
    \begin{equation*}
        \int_{a}^{\infty} f(x)\,dx = \lim_{b \to \infty} \int_{a}^{b} f(x)\,dx.
    \end{equation*}
    if the limit exists. We call this an \textit{improper integral} of $f$. We say the improper integral \textit{converges} if this limit exists and \textit{diverges} otherwise. 
    \begin{enumerate}
        \item 
        Suppose $\int_{a}^{\infty} |f(x)|\,dx$ converges. Prove that $\int_{a}^{\infty} f(x)\,dx$ also converges (this is called \textit{absolute convergence}).\\
        
        \item 
        Suppose further that $a = 1$ and $f(x)$ is nonnegative and monotonically decreasing on $[1, \infty)$. Prove that
        \begin{equation*}
            \int_{1}^{\infty} f(x)\,dx \quad \text{ converges if and only if } \quad \sum_{n = 1}^{\infty} f(n) \quad \text{ converges }.
        \end{equation*}
        This is the \textit{integral test}. \textit{Hint:} draw a picture.\\

        \item 
        Find an example of a function $g : [1, \infty) \to \mathbb{R}$, integrable on $[1, b]$ for all $b > 1$, so that
        \begin{equation*}
            \int_{1}^{\infty} g(x)\,dx
        \end{equation*}
        is convergent but not absolutely convergent. Prove your assertions.
        
    \end{enumerate}

    \vspace{0.5cm}

    \item 
    Let $(X, d_X)$ be a metric space. Suppose $\{f_n\}$ and $\{g_n\}$ are sequences of functions $X \to \mathbb{R}$ that converge uniformly to $f$ and $g$ respectively.
    \begin{enumerate}
        \item 
        Show that the sequence $\{f_n + g_n\}$ converges uniformly to $f + g$.\\

        \item 
        Suppose further that, for every $n \in \mathbb{N}$, the functions $f_n$ and $g_n$ are bounded.\\
        \begin{enumerate}
            \item 
            Show that $f$ and $g$ (and hence $fg$) are bounded.\\
        
            \item 
            Show that the sequences $\{f_n\}$ and $\{g_n\}$ are \textit{uniformly bounded}, in that there is $M > 0$ so that $|f_n(x)| \leq M$ and $|g_n(x)| \leq M$ for all $n \in \mathbb{N}$ and for all $x \in X$.\\

            \item 
            Show that $f_ng_n \to fg$ uniformly.\\

        \end{enumerate}

        \item 
        If we drop the assumption that every $f_n$ is bounded (but keep the assumptions that every $g_n$ is bounded and $f_n \to f$ and $g_n \to g$ uniformly), is it still true that $f_ng_n \to fg$ uniformly? Prove or give a counterexample.
            
    \end{enumerate}

    \vspace{0.5cm}

    \item 
    The purpose of this problem is to define the natural logarithm. You are not allowed to use any properties of this function you might know already (i.e. do not take anything for granted). However, you are allowed to use all of the facts we proved about the exponential function in class.\\
    
    Define a function $f : (0, \infty) \to \mathbb{R}$ via
    \begin{equation*}
        f(x) = \int_{1}^{x} \frac{1}{t}\,dt.\\[5pt]
    \end{equation*}
    \begin{enumerate}
        \item 
        Using the change of variables theorem, prove that
        \begin{equation*}
            f(ab) = f(a) + f(b)
        \end{equation*}
        for all $a, b > 0$.\\

        \item 
        Prove that
        \begin{equation*}
            \lim_{x \to \infty} f(x) = \infty.\\[5pt]
        \end{equation*}

        \item 
        Prove that $f$ is differentiable with
        \begin{equation*}
            f'(x) = \frac{1}{x}
        \end{equation*}
        for all $x \in (0, \infty)$.\\

        \item 
        Prove that $f(e^x) = x$ for all $x \in \mathbb{R}$. \textit{Hint:} differentiate both sides.\\

        \item 
        Prove that $e^{f(x)} = x$ for all $x \in (0, \infty)$ \textit{Hint:} differentiate $\frac{e^{f(x)}}{x}$.\\

        \item 
        Prove that
        \begin{equation*}
            f(x) = \sum_{n = 1}^{\infty} \frac{(-1)^{n + 1}}{n} (x - 1)^n
        \end{equation*}
        for all $x \in (0, 2)$. \textit{Hint:} write $\frac{1}{x}$ as a geometric series.\\

        \item 
        Prove that
        \begin{equation*}
            f(2) = \sum_{n = 1}^{\infty} \frac{(-1)^{n - 1}}{n}.
        \end{equation*}
        \textit{Hint:} you cannot just plug $2$ into part (e) as we do not yet know whether the formula from (e) is valid for $x = 2$ (or whether this series even converges). There is a theorem that will help you here.
        
    \end{enumerate}
    
    \vspace{0.5cm}

    \item 
    For $\alpha \in \mathbb{R}$ and $k \in \{0, 1, 2, \dots\}$, define the \textit{generalized binomial coefficient}
    \begin{equation*}
        \binom{\alpha}{k} = \frac{\alpha(\alpha - 1)(\alpha - 2) \cdots (\alpha - k + 1)}{k!}.
    \end{equation*}
    By convention, we always have $\binom{\alpha}{0} = 1$. The \textit{binomial series} is the power series
    \begin{equation*}
        \sum_{k = 0}^{\infty} \binom{\alpha}{k} x^k.
    \end{equation*}
    \begin{enumerate}
        \item 
        If $\alpha \in \{0, 1, 2, \dots\}$, show that the binomial series converges to $(1 + x)^\alpha$ for all $x \in \mathbb{R}$.\\

        \item 
        If $\alpha \notin \{0, 1, 2, \dots\}$ (which we will assume for the rest of this problem), show that the binomial series has radius of convergence $1$. \textit{Hint:} Ratio test.\\

        \item 
        Let $f : (-1, 1) \to \mathbb{R}$ be defined by $f(x) = \sum_{k = 0}^{\infty} \binom{\alpha}{k} x^k$. Show that $f$ is differentiable and satisfies
        \begin{equation*}
            (1 + x)f'(x)  = \alpha f(x)
        \end{equation*}
        for all $x \in (-1, 1)$.\\
        
        \item 
        Show that $f(x) = (1 + x)^\alpha$ for all $x \in (-1, 1)$ \textit{Hint:} differentiate the function $g(x) = (1 + x)^{-\alpha}f(x)$.\\

        \item 
        Show that 
        \begin{equation*}
            \sqrt{2} = \sum_{k = 0}^{\infty} \binom{1/2}{k}.
        \end{equation*}
        \textit{Hint:}  you cannot just plug in $1$ into the binomial series; the same comments from problem 3 part (f) apply.
        

        
    \end{enumerate}
    

      
    
    
    
    
\end{enumerate}



\end{document}     